% !TeX root = ../main.tex
\section{A Biot coefficient from matched logarithmic stress laws}
\label{sec:stress-trace-biot-construction}

A drained double-prime stress law and a mineral stress law determine the
finite-pressure response through the intrinsic solid stress trace.  We use
the same logarithmic Cauchy-stress form for the skeleton and mineral and derive a
scalar mineral residual suitable for implicit differentiation.  The
construction is elastic, with \(\mbf F^p=\mbf I\) and \(J^e=J\), and supplies the stress-trace construction underlying
Sec.~\ref{sec:numerical-solid-eos}.  At a fixed plastic state, replacing the
independent volume by \(J^e\) gives the same scalar mineral equation; the
conversion to a total-volume tangent is derived in the main text.

\subsection{Skeleton and mineral stress laws}

Prescribe the drained mean double-prime Cauchy stress and the
mean intrinsic mineral stress by
\begin{align}
 \left.\frac13\tr\bs\sigma^{\prime\prime}\right|_{p=0}
 &=\frac{K\ln J}{J},
 \label{eq:trace-drained-skeleton-stress}\\
 \frac13\tr\bar{\bs\sigma}_s
 &=\frac{K_s\ln\bar J}{\bar J}.
 \label{eq:trace-grain-eos-volume-form}
\end{align}
Both laws are derivatives of quadratic logarithmic volume energies with
respect to their corresponding volume ratios.  The factors \(1/J\) and
\(1/\bar J\) therefore place both laws in the same Cauchy-stress form.
Both mean stresses are negative in compression.  The constants \(K\) and \(K_s\) are the reference drained skeleton and mineral
bulk moduli.  We take \(0<K<\phi_{s0}K_s\) for the locally stable reference
state.  The current mineral bulk modulus is
\begin{equation}
 \bar J\frac{\mathrm d}{\mathrm d\bar J}
 \left(\frac13\tr\bar{\bs\sigma}_s\right)
 =\frac{K_s(1-\ln\bar J)}{\bar J}.
 \label{eq:trace-current-mineral-modulus}
\end{equation}
Thus the positive-tangent mineral branch has \(\bar J<\exp(1)\).
The intrinsic solid stress includes the skeleton contribution, while the
density-conjugate pressure \(\bar p_s\) equals pore pressure \(p\) under
local thermodynamic equilibrium.

The volume-weighted phase stresses and the Biot transform give
\begin{align}
 \bs\sigma&=\phi_s\bar{\bs\sigma}_s-(1-\phi_s)p\mbf I,
 \label{eq:trace-phase-stress-decomposition}\\
 \bar{\bs\sigma}_s&=\frac{\bs\sigma^{\prime}}{\phi_s}-p\mbf I,
 \label{eq:trace-intrinsic-solid-stress}\\
 \bs\sigma^{\prime}&=\bs\sigma^{\prime\prime}+(1-B)p\mbf I.
 \label{eq:trace-single-prime-transform}
\end{align}
Taking the trace, as we did in Foster and Xu \citep[App.~B]{fosterxu2025}, yields
\begin{equation}
 \frac13\tr\bar{\bs\sigma}_s
 =\frac{\tfrac13\tr\bs\sigma^{\prime\prime}+(1-B)p}{\phi_s}-p.
 \label{eq:trace-intrinsic-solid-mean-stress-transform}
\end{equation}
Combining the mineral law with this identity and solid mass conservation,
\(\phi_s=\phi_{s0}\bar J/J\), gives
\begin{equation}
 \frac13\tr\bs\sigma^{\prime}
 =\frac{\phi_{s0}}{J}\left(p\bar J+K_s\ln\bar J\right).
 \label{eq:trace-single-prime-mineral-law}
\end{equation}
At zero pore pressure the two effective stresses coincide, so the drained
skeleton law fixes the mineral-volume curve
\begin{equation}
 \bar J(J,0)=J^{K/(\phi_{s0}K_s)}.
 \label{eq:trace-drained-mineral-curve}
\end{equation}
It passes through \(\bar J(1,0)=1\).  Its reference tangent gives
\(B_0=1-K/K_s\); its full form supplies the boundary data away from that
reference point.

\subsection{Finite-pressure closure from stress--volume reciprocity}

Under volumetric variations at fixed isochoric deformation, the pressure
Legendre potential satisfies
\begin{equation}
 \left.\pder{W^{\prime}}{J}\right|_p
 =\frac13\tr\bs\sigma^{\prime},
 \qquad
 \left.\pder{W^{\prime}}{p}\right|_J=\phi_{s0}\bar J.
 \label{eq:trace-legendre-conjugacies}
\end{equation}
Equality of mixed derivatives gives the reciprocal relation
\begin{equation}
 \left.\pder{}{p}\left(\frac13\tr\bs\sigma^{\prime}\right)\right|_J
 =\phi_{s0}\left.\pder{\bar J}{J}\right|_p.
 \label{eq:trace-stress-mineral-reciprocity}
\end{equation}
Substitution of \eqref{eq:trace-single-prime-mineral-law} gives
\begin{equation}
 J\left.\pder{\bar J}{J}\right|_p
 -\left(p+\frac{K_s}{\bar J}\right)
 \left.\pder{\bar J}{p}\right|_J=\bar J.
 \label{eq:trace-compatible-mineral-pde}
\end{equation}
Along its characteristic curves, \(\bar J/J\) and
\(p\bar J+K_s\ln\bar J\) are constant.  On the drained curve,
\(\ln\bar J=[K/(\phi_{s0}K_s)]\ln J\).  Eliminating the drained volume
between the two invariants therefore gives
\begin{equation}
 p\bar J+K_s\ln\bar J
 =\frac{KK_s}{\phi_{s0}K_s-K}\ln\left(\frac{J}{\bar J}\right).
 \label{eq:trace-characteristic-invariant-relation}
\end{equation}
Collecting the logarithmic terms yields the local scalar equation
\begin{equation}
 \boxed{
 K_s\ln\bar J
 +\left(1-\frac{K}{\phi_{s0}K_s}\right)p\bar J
 -\frac{K}{\phi_{s0}}\ln J=0.}
 \label{eq:trace-implicit-mineral-volume}
\end{equation}
Thus finite-pressure equilibrium requires only a scalar constitutive solve.
For \(p\geq0\), the residual is strictly increasing in positive \(\bar J\)
from negative to positive infinity, so the positive root is unique.  For
negative pressure, the branch is continued from the drained state while
retaining a positive local tangent.

\subsection{Implicit state tangent and Biot coefficient}

At fixed pore pressure and referential solid mass, differentiating
\eqref{eq:trace-implicit-mineral-volume} gives
\begin{equation}
 \left[\frac{K_s}{\bar J}
 +\left(1-\frac{K}{\phi_{s0}K_s}\right)p\right]
 \left.\pder{\bar J}{J}\right|_p=\frac{K}{\phi_{s0}J}.
 \label{eq:trace-mineral-tangent-system}
\end{equation}
The resulting tangent and coefficient are
\begin{align}
 \left.\pder{\bar J}{J}\right|_p
 &=\frac{K\bar J}
 {\phi_{s0}J\left[K_s+\left(1-\dfrac{K}{\phi_{s0}K_s}\right)p\bar J\right]},
 \label{eq:trace-mineral-volume-tangent}\\
 B&=1-
 \frac{K\bar J}
 {J\left[K_s+\left(1-\dfrac{K}{\phi_{s0}K_s}\right)p\bar J\right]}.
 \label{eq:trace-closed-form-biot-coefficient}
\end{align}
At \(J=\bar J=1\) and \(p=0\), this gives \(B_0=1-K/K_s\).
On the entire drained path,
\begin{equation}
 B(J,0)=1-\frac{K}{K_s}J^{K/(\phi_{s0}K_s)-1}.
 \label{eq:trace-drained-biot-coefficient}
\end{equation}

The same calculation fits the two-state formulation of
Sec.~\ref{sec:fixed-pressure-constrained-tangent}.  With
\(\mbf y=(\bar\rho_s,\phi_s)\) and
\(\bar J=\bar\rho_{s0}/\bar\rho_s\), use
\begin{align}
 R_m&=J\phi_s\bar\rho_s-J\rho_s=0,
 \label{eq:trace-local-mass-residual}\\
 R_e&=\ln\left(\frac{\bar\rho_s}{\bar\rho_{s0}}\right)
 +\frac{K}{\phi_{s0}K_s}\ln J
 -\left(1-\frac{K}{\phi_{s0}K_s}\right)
 \frac{p\bar\rho_{s0}}{K_s\bar\rho_s}=0.
 \label{eq:trace-local-mineral-residual}
\end{align}
The EOS entries in the local Jacobian and explicit volume derivative are
\begin{equation}
 \pder{R_e}{\bar\rho_s}
 =\frac{1}{\bar\rho_s}\left[
 1+\left(1-\frac{K}{\phi_{s0}K_s}\right)\frac{p\bar J}{K_s}\right],
 \qquad
 \pder{R_e}{J}=\frac{K}{\phi_{s0}K_sJ}.
 \label{eq:trace-local-residual-derivatives}
\end{equation}
The mass-residual derivatives retain their existing form.  Solving
\eqref{eq:implicit-local-state-tangent} therefore recovers
\eqref{eq:trace-mineral-volume-tangent} and
\eqref{eq:trace-closed-form-biot-coefficient}; no differentiation of an
explicit root formula is required.

\subsection{Reconstructed double-prime stress and energy check}

Using the mineral residual to simplify the trace relation gives
\begin{equation}
 \frac13\tr\bs\sigma^{\prime}
 =\frac{K}{J}\left(\ln J+\frac{p\bar J}{K_s}\right).
 \label{eq:trace-closed-form-single-prime-stress}
\end{equation}
The Biot transform then gives the finite-pressure double-prime law,
\begin{equation}
 \frac13\tr\bs\sigma^{\prime\prime}
 =\frac{K}{J}\left\{
 \ln J+
 \frac{\left(1-\dfrac{K}{\phi_{s0}K_s}\right)(p\bar J)^2}
 {K_s\left[K_s+\left(1-\dfrac{K}{\phi_{s0}K_s}\right)p\bar J\right]}
 \right\}.
 \label{eq:trace-closed-form-double-prime-stress}
\end{equation}
The pressure correction is quadratic near the reference state and vanishes
on the entire drained path.  It follows from the mineral law and reciprocity.
For a decoupled neo-Hookean isochoric response, add
\(GJ^{-5/3}\operatorname{dev}(\mbf F\mbf F^T)\) to this mean stress times
\(\mbf I\).  The total stress is
\(\bs\sigma=\bs\sigma^{\prime\prime}-Bp\mbf I\).

A compact verification uses the reduced energy evaluated on the mineral state,
\begin{equation}
 W^{\prime\prime}(\mbf F,p)
 =\frac{G}{2}\left[J^{-2/3}\tr(\mbf F\mbf F^T)-3\right]
 +\frac{K}{2}(\ln J)^2
 +\frac{\phi_{s0}}{2K_s}
 \left(1-\frac{K}{\phi_{s0}K_s}\right)(p\bar J)^2.
 \label{eq:trace-compatible-reduced-potential}
\end{equation}
Its fixed-pressure deformation derivative gives the double-prime stress
above.  Its pressure derivative is
\(-\phi_{s0}p\,\partial\bar J/\partial p\).  Adding
\(\phi_{s0}p\bar J\) gives \(W^{\prime}\) and recovers both conjugacies in
\eqref{eq:trace-legendre-conjugacies}.  These checks retain the implicit
dependence of \(\bar J\) on \(J\) and \(p\).

\subsection{Reference response and admissible branch}

The pressure tangent follows from the same scalar residual,
\begin{equation}
 \left.\pder{\bar J}{p}\right|_J
 =-\frac{\left(1-\dfrac{K}{\phi_{s0}K_s}\right)\bar J^2}
 {K_s+\left(1-\dfrac{K}{\phi_{s0}K_s}\right)p\bar J}.
 \label{eq:trace-mineral-pressure-tangent}
\end{equation}
At reference it gives
\begin{equation}
 \left.\pder{(J-\phi_{s0}\bar J)}{p}\right|_0
 =\frac{\phi_{s0}}{K_s}-\frac{K}{K_s^2}.
 \label{eq:trace-reference-pore-volume-storage}
\end{equation}
Adding fluid compressibility recovers
\eqref{eq:classical-biot-storage-comparison}.
Under unjacketed hydrostatic loading, \(J=\bar J\),
\(p=-K_s\ln J/J\), and \(\phi_s=\phi_{s0}\).  The single-prime stress
vanishes and the total stress is \(-p\mbf I\).  In the incompressible-mineral
limit \(K_s\to\infty\), \(\bar J\to1\), \(B\to1\), and the
double-prime law reduces to the drained skeleton response.

Admissible states require \(J>0\), \(0<\bar J<\exp(1)\), and
\(0<\phi_{s0}\bar J/J<1\).  A stable local mineral equilibrium requires
\begin{equation}
 K_s+\left(1-\frac{K}{\phi_{s0}K_s}\right)p\bar J>0.
 \label{eq:trace-mineral-stability-domain}
\end{equation}
This condition makes the density Jacobian nonsingular and the mineral
pressure tangent negative.  At negative pore pressure it selects the branch
connected to the drained state; it does not assert global existence or
stability.  The fixed-plastic-state extension and its coupled numerical
implementation are given in Secs.~\ref{sec:inelastic-plastic-driving-stress}
and~\ref{sec:implicit-ad-implementation}.
